\section{ماژول erc\_svd (Error-Controlled SVD Compressor)}
\label{sec:module_erc_svd}

\subsection{معرفی و انگیزه}
این مکانیزم بر مبنای مبانی نظری پیشرفته در ادبیات توجه و فشرده‌سازی \cite{ref_ercsvd_2024} توسعه یافته است.

در بسیاری از سناریوهای فشرده‌سازی لایه‌های ترنسفورمر، تحمیل یک رتبه ثابت از پیش تعیین‌شده به همه لایه‌ها غیرواقع‌بینانه است؛ چرا که لایه‌های مختلف دارای زوال طیفی متفاوتی در مقادیر منفرد خود هستند. برخی لایه‌ها با نگه‌داشتن تعداد اندکی از مقادیر منفرد بیش از ۹۵ درصد از انرژی وزن را حفظ می‌کنند، در حالی که لایه‌های دیگر نیاز به رتبه بالاتری برای حفظ کارکرد خود دارند. در نتیجه، تنظیم مبتنی بر رتبه اغلب یا به فشرده‌سازی ناکافی یا به تخریب غیرقابل جبران دقت می‌انجامد.

ماژول فشرده‌ساز \lr{ERC-SVD} با نام ثبتی \lr{erc} و \lr{erc\_svd} («تجزیه مقادیر منفرد با خطای کنترل‌شده» یا \lr{Error-Controlled SVD})، با اتخاذ رویکرد دوگان، هدف فشرده‌سازی را بر پایه یک سقف خطای نسبی فروبنیوس مجاز (\lr{Bounded Frobenius Truncation Error}) تعریف می‌کند. در این راهبرد، به ازای هر لایه، کمترین رتبه ممکن $r$ به گونه‌ای انتخاب می‌شود که تضمین کند خطای برش مقادیر منفرد از سقف مشخص‌شده فراتر نرود. با این کار، توازن بهینه‌ای میان نرخ تراکم و پایداری عملکرد مدل برقرار می‌گردد.

\subsection{مبانی نظری و فرمول‌بندی ریاضی}
برای ماتریس وزن $\mathbf{W} \in \mathbb{R}^{d_{\text{out}} \times d_{\text{in}}}$ با تجزیه مقادیر منفرد $\mathbf{W} = \mathbf{U} \mathbf{\Sigma} \mathbf{V}^\top$، مقادیر منفرد به ترتیب نزولی $\sigma_1 \ge \sigma_2 \ge \dots \ge \sigma_d \ge 0$ مرتب شده‌اند. انرژی کل ماتریس بر حسب نُرم فروبنیوس برابر است با:
\begin{equation}
\|\mathbf{W}\|_F^2 = \sum_{i=1}^d \sigma_i^2
\end{equation}
تقریب رتبه $r$ ماتریس وزن به صورت $\mathbf{W}_r = \sum_{i=1}^r \sigma_i \mathbf{u}_i \mathbf{v}_i^\top$ دارای خطای برشی با فرم بسته زیر است:
\begin{equation}
\|\mathbf{W} - \mathbf{W}_r\|_F^2 = \sum_{i=r+1}^d \sigma_i^2
\end{equation}
با تعیین آستانه خطای مجاز $\epsilon_{\text{allowed}} \in (0, 1)$، ماژول کمترین رتبه $r^*$ را بر اساس نسبت انرژی انباشته محاسبه می‌کند:
\begin{equation}
\frac{\sum_{i=r^*+1}^d \sigma_i^2}{\sum_{i=1}^d \sigma_i^2} \le \epsilon_{\text{allowed}}^2 \iff \frac{\sum_{i=1}^{r^*} \sigma_i^2}{\sum_{i=1}^d \sigma_i^2} \ge 1 - \epsilon_{\text{allowed}}^2
\end{equation}
با به دست آمدن $r^*$، با شرط عدم تجاوز از سقف بازدهی فشرده‌سازی $\frac{d_{\text{in}} d_{\text{out}}}{d_{\text{in}} + d_{\text{out}}}$، فاکتورهای چپ و راست با تفکیک متقارن توان مقادیر منفرد ساخته می‌شوند:
\begin{equation}
\mathbf{U}_{r^*} = \mathbf{U}_{:, :r^*} \mathbf{\Sigma}_{:r^*}^{1/2}, \quad \mathbf{V}_{r^*} = \mathbf{\Sigma}_{:r^*}^{1/2} \mathbf{V}_{:, :r^*}^\top
\end{equation}
و لایه با ماژول فاکتورگیری‌شده جایگزین می‌گردد:
\begin{equation}
\mathbf{W}_{\text{comp}} = \mathbf{U}_{r^*} \mathbf{V}_{r^*}, \quad \mathbf{y} = (\mathbf{x} \mathbf{V}_{r^*}^\top) \mathbf{U}_{r^*}^\top + \mathbf{b}
\end{equation}

\paragraph{تحلیل طیفی و نامساوی‌های پایداری خطا}
برای بررسی میزان دقت تقریب ماتریس وزن یا نگاشت توجه در ماژول \lr{erc\_svd}، تفاضل ماتریس فاکتورگیری‌شده از نگاشت متراکم اصلی بر اساس نامساوی هافدینگ-فون‌نیومان ارزیابی می‌گردد. فرض کنید $\mathbf{W} \in \mathbb{R}^{d_{\text{in}} \times d_{\text{out}}}$ ماتریس وزن هدف باشد که توسط تقریب کم‌مرتبه یا پراکنده با رتبه موثر $r$ جایگزین شده است. کران بالای خطای بازسازی بر مبنای نرم فروبنیوس به صورت زیر تصریح می‌شود:
\begin{equation}
\label{eq:spectral_bound_erc_svd}
\| \mathbf{W} - \mathbf{W}_{r} \|_F = \sqrt{\sum_{i=r+1}^{\min(d_{\text{in}}, d_{\text{out}})} \sigma_i^2(\mathbf{W})} \le \epsilon \| \mathbf{W} \|_F
\end{equation}
که در آن $\sigma_i(\mathbf{W})$ بیانگر $i$-امین مقدار منفرد مرتب‌شده ماتریس به صورت نزولی است. نسبت حفظ انرژی طیفی $\eta_r = \frac{\sum_{i=1}^r \sigma_i^2}{\sum_{i=1}^{d} \sigma_i^2}$ تضمین می‌کند که با انتخاب آستانه $\tau \ge 0.90$، بیش از نود درصد از پراکندگی واریانس داده‌ها حفظ شده و افت سیگنال در انتشار گرادیان رو به عقب به حداقل ممکن کاهش می‌یابد.

\paragraph{تحلیل مجانبی و پیچیدگی اشتراک حافظه سخت‌افزاری}
از منظر پیچیدگی محاسباتی و انتقال داده‌ها در سلسومراتب حافظه (حافظه سراسری \lr{HBM} به حافظه محلی \lr{SRAM})، عملیات ضرب ماتریسی در این مکانیزم نیازمند تخصیص بافر موقت با ابعاد فاکتورگیری‌شده است. برای پردازش دنباله‌ای به طول $N$ با بعد مخفی $D$، بار محاسباتی هسته از فرم درجه دو به رابطه زیر تقلیل می‌یابد:
\begin{equation}
\label{eq:complexity_erc_svd}
\mathcal{C}_{\text{erc\_svd}} = \mathcal{O}\left( B \cdot N \cdot D \cdot r + B \cdot H \cdot N \cdot d_k \right) \ll \mathcal{O}\left( B \cdot H \cdot N^2 \cdot d_k \right)
\end{equation}
این بهبود مرتبه پیچیدگی، با کاهش چشمگیر حجم تبادل تانسورها در گذرگاه‌های ارتباطی کارت گرافیک، مانع از اشباع پهنای باند حافظه (پدیده \lr{Memory-Bound}) در فرآیند استنتاج طولانی‌مدت می‌گردد.

\subsection{معماری و پیاده‌سازی}
پیاده‌سازی کلاس \lr{ERCCompressor} در فایل \lr{Codes/framework/compression\_engine/erc\_svd.py} انجام شده و شامل ویژگی‌های مهندسی زیر است:
\begin{itemize}
    \item \textbf{آزمون انباشت انرژی تجمعی (\lr{Cumulative Energy Search}):} استفاده از \lr{torch.cumsum} بر روی توان دوم مقادیر منفرد برای تعیین سریع و برداری شاخص $r^*$ بدون جستجوی ترتیبی پرهزینه.
    \item \textbf{سازگاری کامل با \lr{PEFT} و آداپتورهای \lr{LoRA}:} شناسایی رتبه \lr{LoRA} از طریق بعد تانسور \lr{lora\_A.weight} و انتقال ضریب مقیاس‌گذاری (\lr{lora\_alpha}) به لایه فشرده جدید.
    \item \textbf{آستانه بازگشت و ایمنی (\lr{Fallback Guard}):} در صورتی که پس از تجزیه، خطای عملی فروبنیوس به دلیل خطای محاسباتی ممیز شناور از \lr{max\_error} فراتر رود، لایه اصلی بدون دستکاری بازگردانده شده و انجماد پارامتر اعمال می‌شود.
    \item \textbf{پشتیبانی از لایه‌های ترنسفورمرهای متنی متنوع:} سازگاری با لایه‌های خطی استاندارد در انکودرها و لایه‌های پروجکشن \lr{Conv1D} در ترنسفورمرهای خودبازگشتی.
\end{itemize}

\subsection{نتایج تجربی}
نتایج ارزیابی تجربی ماژول \lr{erc\_svd} بر روی مدل‌های مرجع در جدول \ref{tab:benchmark_erc_svd} ارائه شده است:

\begin{table}[H]
\centering
\small
\caption{نتایج بنچ‌مارک تجربی فشرده‌ساز \lr{erc\_svd} بر روی مدل‌های پایه}
\label{tab:benchmark_erc_svd}
\begin{tabular}{lccccc}
\toprule
\textbf{مدل پایه \lr{(Base Model)}} & \textbf{دقت \lr{(Accuracy)}} & \textbf{شاخص \lr{F1}} & \textbf{حافظه \lr{(Peak VRAM)}} & \textbf{توان \lr{(Throughput)}} & \textbf{وضعیت} \\
\midrule
\lr{DistilBERT} & \lr{90.71\%} & \lr{0.9103} & \lr{418.17 MB} & \lr{236,111.95 tok/s} & \lr{PASS} \\
\lr{RoBERTa} & \lr{94.15\%} & \lr{0.9425} & \lr{503.28 MB} & \lr{14,452.28 tok/s} & \lr{PASS} \\
\lr{GPT-2} & -- & -- & -- & -- & \lr{NOT\_RUN} \\
\lr{Qwen} & -- & -- & -- & -- & \lr{SKIP} \\
\bottomrule
\end{tabular}
\end{table}

بر اساس اندازه‌گیری‌های آزمایشگاهی، کنترل خطا با تضمین نظری موجب شده مدل \lr{DistilBERT} به دقت چشمگیر \lr{90.71\%} و توان قابل توجه \lr{236,111.95 tok/s} دست یابد. در مدل \lr{RoBERTa} نیز دقت \lr{94.15\%} با شاخص F1 برابر \lr{0.9425} و مصرف حافظه بهینه \lr{503.28 MB} ثبت گردیده است. اجرای ماژول برای \lr{GPT-2} ثبت نگردیده (\lr{NOT\_RUN}) و برای \lr{Qwen} در وضعیت \lr{SKIP} قرار دارد.

\subsection{محدودیت‌ها و تحلیل شکست}
اگرچه کنترل خطای فروبنیوس از جنبه ریاضی خوش‌رفتار است، اما نُرم فروبنیوس اهمیت نسبی سطرهای مختلف یا جهات ورودی فعال‌سازی را در نظر نمی‌گیرد؛ یعنی خطای برش در جهت‌هایی از وزن که در لایه‌های پیشین ورودی کمی دریافت می‌کنند با خطای در جهت‌های پرتراکم به صورت یکسان جریمه می‌شود. در نتیجه در برخی ساختارها، برای رسیدن به خطای پایین، رتبه بالایی تخصیص می‌یابد که ممکن است نرخ فشرده‌سازی نهایی را محدود کند.
